Els dos satèl·lits de Mart, Fobos i Deimos, porten el nom dels fills bessons d'Afrodita i Ares. En la mitologia romana, Ares, déu de la guerra, s'identifica amb Mart. Els dos satèl·lits tenen forma irregular, però els podem aproximar a una esfera de diàmetre 22,2~km per a Fobos i de 12,6~km per a Deimos. Per tant, comparats amb la Lluna, que té un diàmetre de 3\,475~km, són petits. El radi orbital mitjà (distància entre els centres dels dos objectes) de Fobos al voltant de Mart és de 9\,377~km i el seu període de revolució és de 7~hores, 39~minuts i 14~segons. Sabent que el radi orbital mitjà de Deimos és de 23\,460~km, determineu a partir d'aquestes dades:

\begin{apartat}{1,25}
La massa de Mart i la intensitat del camp gravitatori que Mart crea a la seva superfície.
\end{apartat}

\begin{apartat}{1,25}
El període de revolució de Deimos al voltant de Mart i la seva energia mecànica.
\end{apartat}

\dades{%
  $G = 6,67\times 10^{-11}\ \mathrm{N\,m^2\,kg^{-2}}$.\\
  $M_\mathrm{Fobos} = 1,10\times 10^{16}\ \mathrm{kg}$.\\
  $M_\mathrm{Deimos} = 2,00\times 10^{15}\ \mathrm{kg}$.\\
  $R_\mathrm{Mart} = 3\,390\ \mathrm{km}$.}
\nota{Considereu que els dos satèl·lits descriuen una trajectòria circular al voltant de Mart.}
